\documentclass{report}
\usepackage{amsmath,amssymb}
\setlength{\parindent}{0mm}
\setlength{\parskip}{1em}
\begin{document}
\begin{center}
\rule{6in}{1pt} \
{\large Chad Westphal \\
{\bf FOSLL$^*$ For Nonlinear Partial Differential Equations}}

Dept of Math \& CS \\ Wabash College \\ Crawfordsville \\ IN 47933
\\
{\tt westphac@wabash.edu}\\
Eunjung Lee\\
Thomas Manteuffel\end{center}

Least-squares finite element methods are designed around the idea of
minimizing discretization error in an appropriate norm. For sufficiently
regular, elliptic-like problems, a first-order system least squares
(FOSLS) approach may be continuous and coercive in the $H^1$ norm,
yielding solutions accurate in $H^1$. For problems posed in high aspect
ratio domains (e.g., flow through a long channel), approximations on
coarse resolutions may have error whose $L^2$ norm is large relative to
the $H^1$ seminorm. The first order system $LL^*$ (FOSLL$^*$) approach
seeks to minimize the residual of the equations in a dual norm induced by
the differential operator, yielding a better $L^2$ approximation. Mass
conservation in fluid flow, for example, is greatly enhanced by such an
approach. In this talk, we extend this general framework to nonlinear
problems.

Newton's method is a typical outer iteration for an efficient finite
element approximation of nonlinear partial differential equations. We
present the framework for an inexact Newton iteration based on a
FOSLL$^*$ approximation to each linearization step and establish theory
for convergence. Numerical results are presented for a
velocity-vorticity-pressure formulation of the steady incompressible
Navier-Stokes equations, and we discuss extensions and comparisons to the
more typical Newton-FOSLS approach.


\end{document}
